\chapter{SMT Representation of Policies}
\label{ch:04}

% In this section
%% Why we transform to SMT
%% How will this chapter be structured
%%% Transformation of incrementally more complex constructs
As the analysis of Rego policies leverages the power of SMT solving, it is crucial to have a function for transforming policy constructs into corresponding SMT statemenents.
This section will provide theoretical information about this process.
Due to the complexity and expressiveness of the language, the description will start with covering simple scenarios and gradually introduce new constructs for this procedure to support, according to the structure of chapter \ref{ch:form}.
It is important to note that not all Rego statemenents are fully representable in SMT, which we will discuss in the last section of this chapter (\todo{re}).

\section{Variables and Literals}

Since variables and literals stand in the core of all policies, we will first look at their representation and reason about its effectiveness.

\subsection{Atomic values}

Rego allows for atomic values of these types:
\begin{itemize}
  \item String
  \item Number - natural numbers (32 bit) and decimals (float by IEEE 754)
  \item Boolean - either true or false
  \item Null - special type for the \texttt{null} value
\end{itemize}

From these types, both Strings and Booleans are strictly representable in SMT-LIB2 and Null type can be simply substituted by a custom type holding only one value.
As it was described in section \ref{sec:smt}, SMT-LIB2 supports both BitVector and Float sorts, which could be unionized with a custom type and this directly match the Rego Number type.

On the other hand, the set of policy examples \todo{cite} provided by the Rego reference does not include any decimal values, and thus my supervisor for this work and I decided, that support for this type is expendable.
% Furthermore, the direct BitVector representation brings up 
Furthermore, with the policy as code principle stresses the readability of rules, it is expected for them not to rely on obscure representation-based behaviours such as bit overflow.
For this cause, we decided that translating natural number in Rego into the SMT-LIB2 Int values is also sufficient.
% and also easily interchangeable with the BitVector kind.
For the rest of this paper, we will thus use the transformation of the Rego Number type into the SMT-LIB2 Int, as it is more clear to understand than the latter option.

% maybe add some sort of table which shows, how different values are mapped.
% for example 1 -> 1, "ahoj" -> "ahoj", null -> ONull
With Rego not having a strong type system, which allows variables to be for example either Int or String, we will (for now) encode their type in SMT as a unionizing sort over all atomic possibilities.
This brings us to the sort \texttt{OTypeD0} defined as:

% TODO: maybe rename OTypeD0 as OTypeAtom
\begin{lstlisting}
(declare-datatypes ()
  ((OTypeD0
    (OString (str String))
    (ONumber (num Int))
    (OBoolean (bool Boolean))
    ONull
  )))
\end{lstlisting}

For custom defined SMT datatypes we use the prefix "O".
The suffix "D0" implies the depth equal to $0$, whose reason will become more clear with the support of more complex types.
With variables being of this sort, we can assign them any of the said atomic types.

\subsection{Objects}

As OPA operates primarily on key-value collections (as described in section \ref{sec:opa}), the SMT support for objects becomes the most crucial part of the whole analysis.
The main problem with objects is that they can hold values, that are also objects, which causes the depth to be unbounded.
Such characteristic can be represented with either a recursive datatype, or set of datatypes of various depths, going up to the maximal depth appearing in the input policy.
In this work, we will focus mainly on the latter option, since it is more easily transformable into the theory of uninterpreted function, which SMT solvers can efectively reason about.

% Rego objects != JSON objects
%% in rego: keys can be of any type
%% in JSON: keys only strings

\subsubsection{Unnested objects}

As mentioned in \todo{ref}, SMT-LIB2 offers an Array sort, which maps keys of one kind to values of another kind.
If we only considered unnested Rego objects, such as
\begin{equation}
  \label{03:eq:simple-object}
  \{ "a": 1, "b": true \}
\end{equation}
we could possibly interpret them in SMT with construct
\begin{center}
  \texttt{Array OTypeD0 OTypeD0}
\end{center}

In adition, OPA communicates with JSON objects, whose keys can only be Strings.
This creates an inconsistency between the inner representation of values in Rego and its JSON output, that is very treaturous to deal with using a formal verification tool.
For this, we allow objects to be represented with more simplified sort
\begin{center}
  \texttt{Array String OTypeD0}
\end{center}

The problem with this solution in our current SMT representation is that Arrays are total mappings, which means that every possible key has a value.
This property requires that we need to explicitly specify for some keys, that they do not have a corresponding value, which is equivalent to assigning them a value of a special sort, which will be used only for undefined values.
With this need in mind we will update our current SMT sort for atomic values:

\begin{lstlisting}[language=lisp]
(declare-datatypes ()
  ((OTypeD0
    (OString (str String))
    (ONumber (num Int))
    (OBoolean (bool Bool))
    ONull
    OUndef)))
\end{lstlisting}

We can now leverage constant operators and store operations provided by state-of-the-art SMT solvers, such as \ziii, to correctly represent on object with only certain keys set.
Going back to the simple example \ref{03:eq:simple-object}, its transformation into SMT could yield a representation such as:
\begin{lstlisting}[language=lisp]
(store
  (store
    (as-const (Array OTypeD0 OTypeD0) OUndef)
    "a" (ONumber 1))
  "b" (OBoolean true))
\end{lstlisting}

  This construct firstly creates an empty object with the \texttt{as-const} expression, i.e., an object whose keys all map to value \texttt{OUndef}.
  The \texttt{store} operation firstly assign this empty object key "a" mapping to Integer value 1, and then, with the top-level function, mapping key "b" to the Boolean value \texttt{true}.

\subsubsection{Nested objects}

We already mentioned, that this work will prefer specifying a separate sort for each possible depth of objects over using a recursive declaration.
Such sort of depth $N$ has to hold at least these two requirements:
  \begin{enumerate}
    \item to have its keys hold value of lower depth,
    \item to be tested, whether it is undefined.
  \end{enumerate}
For the first requirement, a small problem arises from the fact, that elements of an object do not explicitly share the same depth, for example with this object such as:
\begin{center}
  \{ "D0": 1, "D1": \{ "foo": "bar" \} \}.
\end{center}
If we were to enumerate all the lower depths of every new kind, the output SMT representation might simply be populated with loads of excess declarations.
  Other option is to allow for \texttt{OTypeD<N>} to be either an object of values with kind \texttt{OTypeD<N-1>}, or to be interpreted strictly as a value of this kind.
  This way, the datatype declaration would consist of at least two interpretations, described with the following SMT:

\begin{lstlisting}[language=lisp]
(declare-datatypes ()
  ((OTypeD<N>
    (Wrap<N> (wrap<N> OTypeD<N-1>))
    (OObj<N> (obj<N> (Array String OTypeD<N-1>)))
  )))
\end{lstlisting}

  For the second requirement, we can either create a level-specific undefined values (i.e., \texttt{OUndefD2}), or have a way to interpret \texttt{OTypeD<N>} as the atomic \texttt{OTypeD0} which has this feature.
  Compared by the length of code added to the kind declaration, both of the suggested solution are equal.
  For testing whether a certain value is undefined, the level-specific OUndef declaration yields smaller output SMT, since the other option needs to firstly specify that the value is interpreted as atomic and then test it.
  Even still, having unified kind for undefined values allows for quantified tests over object elements.
  For example, checking emptiness of object $o$ of elements $e$ of depth $N$ could be translated as evaluating the following predicate:
  $$
  \neg\exists_{e \in o}: (\texttt{is-OUndef} \hspace{1mm} (\texttt{Atom}_N \hspace{1mm} e))
  $$
  
  This statement is much more complicated in objects with level-specific undefined values, whose elements have different depth.
  For example, let $O_3$ be an object of depth $3$, and $O_1$ an object of depth $1$.
  For Rego expression
  \begin{center}
    $O3$["a"] := $O1$
  \end{center}
  our current translating procedure would yield the following SMT code:
  \begin{center}
    \texttt{(assert (select O3 "a" (Wrap2 O1)))}
  \end{center}
  The important part in this demonstration is that the value of $O_1$ stored in $O_3$ is wrapped in constructor for kind \texttt{OTypeD2}.
  So the naively modified predicate for checking emptiness of object
  $$
  \neg\exists_{e \in o}: (\texttt{is-OUndef2} \hspace{1mm} e)
  $$
  would always be unsatisfied even if $O_1$ had value of type \texttt{OUndef1}.
  To make this predicate correct, we would have to take into account every possible depth of element in the tested object.

% isnt the option of multiple level-specific undefined values better?
% level-specific undefined
%% + smaller when asking for specific values
%% - when object = { "a": a, "b": b }, and a,b have different depths, how do we determine, that this object is empty?


  In summary, compared to the unified undefined kind, level-specific undefined values are slightly more efficient in checking specific values, but much more complex to check over quantified object elements.
  For the rest of this paper, we will continue with using only the atomic undefined type, and thus we extend our representation of variable types of depth $N$ with a special option to extract the atomic value straight directly:
\begin{lstlisting}[language=lisp]
(declare-datatypes ()
  ((OTypeD<N>
    (Atom<N> (atom<N> OTypeD0))
    (Wrap<N> (wrap<N> OTypeD<N-1>))
    (OObj<N> (obj<N> (Array String OTypeD<N-1>)))
  )))
\end{lstlisting}
  
  \subsection{Arrays}

  Rego supports heterogenous arrays, meaning its values do not have to hold the same type.
  For this reason, to translate them into SMT, it is also necessary to count the maximal depth of elements of given array.
  Like objects, arrays are collections mapping keys to values, with the exception that its keys have to be natural numbers starting from 0 and gradually increasing.
  We can thus represent them, for elements of depth $N$, as
  \begin{center}
    \texttt{Array Int OTypeD<N>}
  \end{center}
  with adding mentioned constraints on the keys, so that given models do not have unreasonable array indexing such as with negative numbers.

  As mentioned in section \todo{\ref{sec:smt}} modern SMT solvers, such as \ziii or \todo{cvc5}, also add support for special sort \texttt{Seq}, which has these constraints built in, and thus is more suitable for representing Rego arrays.
  Representing a variable type of depth $N$ will with the support for Rego arrays look as the following:
\begin{lstlisting}[language=lisp]
(declare-datatypes ()
  ((OTypeD<N>
    (Atom<N> (atom<N> OTypeD0))
    (Wrap<N> (wrap<N> OTypeD<N-1>))
    (OObj<N> (obj<N> (Array String OTypeD<N-1>)))
    (OArray<N> (arr<N> (Seq OTypeD<N-1>)))
  )))
\end{lstlisting}


  \subsection{Sets}

  In Rego, sets are also heterogenous collections, but their behaviour is equivalent to the SMT \texttt{Set} sort, provided by the state-of-the-art solvers.
  The only important thing to note is the incompatibility of sets with JSON documents, which do not support them.
  For this cause, after evaluating the policy, all sets have to be transformed into JSON arrays.
  Other than that, we can simply add them to the current definition of variable types:
\begin{lstlisting}[language=lisp]
(declare-datatypes ()
  ((OTypeD<N>
    (Atom<N> (atom<N> OTypeD0))
    (Wrap<N> (wrap<N> OTypeD<N-1>))
    (OObj<N> (obj<N> (Array String OTypeD<N-1>)))
    (OArray<N> (arr<N> (Seq OTypeD<N-1>)))
    (OSet<N> (set<N> (Set OTypeD<N-1>)))
  )))
\end{lstlisting}

\section{Expressions}

% How can expressions look
% references
% basic operations
% builtin functions
%% mention assigning calls - plus(1,2,output)
% local variables
% special keywords
%% not
%% quantifiers - some, every
%% contains 

Previously estabilished conversion of literals and variables from Rego policy language into SMT enables transformation of expessions.
Since Rego emphasises on expressiveness, it provides handful of different ways to form such expressions.

\subsection{References}

References are expressions used to access Rego collections, such as objects, arrays, or sets.
They start with the given collection, followed by a chain of squared brackets, enclosing accessed keys, e.g.:
\begin{center}
  \texttt{obj["a"][0]["b"]}
\end{center}
If the key is a string literal, the squared brackets can be replaced by a dot followed by the literal value, resembling syntax of modern programming languages (e.g., Python, or Javascript).
So the above example can be equivalently written as
\begin{equation}
  \label{04:eq-reference}
  \texttt{obj.a[0].b}
\end{equation}

% TODO: define type analysis beforehand
To translate a reference we use the previously defined function $\TA$ for getting types of terms.
Starting with translation of the base collection (\texttt{obj} in the given example) and extracting its collection type value from the variable sort (with datatype members as \texttt{obj<N>}, \texttt{arr<N>}, or \texttt{set<N>}), we incrementally wrap this SMT value with selectors for each following value in the chain.
Transforming the Rego from example \ref{04:eq-reference} could yield (based on the specific types of given collections):
\begin{lstlisting}[language=SMT]
(select 
  (obj1 (seq.nth
    (arr2 (select 
      (obj3 obj)  ; base object obj
      "a"))       ; obj.a
    0))           ; obj.a[0]
  "b")            ; obj.a[0].b
\end{lstlisting}

For a collection with depth $N$, it is also important to check the actual depth of the selected term as it does not have to equal the expected depth $N-1$.
In this case, we need to access the lower depth interpretation with possibly chained \texttt{wrap} operators.
Another possible transformation of the current example could be:
\begin{lstlisting}[language=SMT]
(select 
  (obj1 (seq.nth
    (arr2 (wrap3 (select  ; accessing lower depth representation
      (obj4 obj)
      "a"))
    0)))
  "b")
\end{lstlisting}
This code would be generated for $|\TA(obj)| = 4$ and $|\TA(obj.a)| = 2$.





